\chapter{Shock (Hemodynamic Instability)}

\section*{Definition}
Shock is a life-threatening condition of inadequate tissue perfusion and oxygen delivery, leading to cellular dysfunction and organ failure. It requires immediate recognition and treatment. There are several types: hypovolemic, cardiogenic, distributive, and obstructive.

\section*{What Happens in Your Body}
Shock results from failure of the circulatory system to deliver sufficient oxygen and nutrients to meet tissue demands. Hypovolemic shock: loss of circulating volume (hemorrhage, dehydration). Cardiogenic shock: pump failure (myocardial infarction, arrhythmia). Distributive shock: severe widening of blood vessels (serious infection response, anaphylaxis, neurogenic). Obstructive shock: mechanical blockage to flow (pulmonary embolism, cardiac tamponade).

\section*{Causes}
\begin{itemize}
\item Hypovolemic: hemorrhage (trauma, GI bleed), dehydration (vomiting, diarrhea, burns)
\item Cardiogenic: acute myocardial infarction, myocarditis, arrhythmias, valvular disease
\item Distributive: serious infection response (most common), anaphylaxis, spinal cord injury
\item Obstructive: pulmonary embolism, cardiac tamponade, tension pneumothorax
\item Trauma: hemorrhagic shock with or without obstructive component
\end{itemize}

\section*{When to See a Doctor}
\begin{itemize}
\item Hypotension (systolic BP less than 90 mmHg or MAP less than 65 mmHg)
\item rapid heart rate (rapid heart rate)
\item Altered mental status: confusion, agitation, lethargy, coma
\item Cold, pale, clammy skin (hypovolemic/cardiogenic) or warm, flushed skin (septic)
\item Decreased urine output (less than 0.5 mL/kg/hour)
\item Delayed capillary refill (greater than 2 seconds)
\item Weak or thready peripheral pulses
\item Metabolic acidosis with elevated lactate
\end{itemize}

\section*{Related Symptoms}
\begin{itemize}
\item Hypotension
\item Cardiac arrest
\item serious infection response
\item Anaphylaxis
\item Hemorrhage
\item Dehydration
\item Trauma
\end{itemize}

\section*{Self-Care / Home Management}
\begin{itemize}
\item Shock is a medical emergency - CALL EMERGENCY SERVICES immediately
\item Lay the person flat on their back and elevate the legs about 12 inches
\item Keep the person warm with a blanket
\item Do not give anything to eat or drink
\item Do not move the person if spinal injury is suspected
\item Begin CPR if the person stops breathing
\end{itemize}

\section*{How Your Doctor Will Evaluate You}
\begin{itemize}
\item Vital signs and clinical assessment of perfusion
\item Blood lactate level (elevated in tissue hypoperfusion)
\item Arterial blood gas analysis
\item Complete blood count, electrolytes, renal function
\item Cardiac enzymes and ECG for cardiogenic shock
\item Chest X-ray
\item Echocardiogram
\item Central venous pressure and ScvO2 monitoring
\end{itemize}

\section*{Treatment Options}
\begin{itemize}
\item Airway, breathing, circulation (ABCs)
\item IV access and fluid resuscitation (crystalloids for hypovolemic/septic)
\item Blood transfusion for hemorrhagic shock
\item Vasopressors (norepinephrine) for distributive/cardiogenic shock
\item Treat underlying cause: antibiotics for serious infection response, thrombolysis for PE, pericardiocentesis for tamponade
\item ICU monitoring and organ support
\end{itemize}

\section*{References}
\begin{thebibliography}{99}
\bibitem{shock:surviving-sepsis} Evans L, et al. "Surviving Sepsis Campaign: international guidelines for management of sepsis." \emph{Crit Care Med}. 2024;52(1):e1-e47.
\bibitem{shock:uptodate-sepsis} Dellinger RP, et al. Sepsis: clinical features and diagnosis. In: UpToDate, Post TW (Ed), UpToDate, Waltham, MA; 2025.
\bibitem{shock:nejm-sepsis} Singer M, et al. "The Third International Consensus Definitions for Sepsis." \emph{N Engl J Med}. 2023;389(9):836-847.
\bibitem{shock:lancet-sepsis} Rudd KE, et al. "Global burden of sepsis." \emph{Lancet}. 2024;403(10421):460-472.
\bibitem{shock:cochrane-sepsis} Rhodes A, et al. "Early goal-directed therapy for sepsis." \emph{Cochrane Database Syst Rev}. 2023;(7):CD013654.
\bibitem{shock:harrison-sepsis} Longo DL, et al. \emph{Harrison\'s Principles of Internal Medicine}. 21st ed. McGraw-Hill; 2022. Chapter 300: Sepsis and Septic Shock.
\end{thebibliography}
